\chapter{PYQs: Colonial Knowledge and the Faction-Ridden Village}

\section{PYQ 2020: Colonial Administrators and the Construction of Traditionalism}

\begin{PYQLink}
\begin{itemize}
\item \textbf{PYQ (2020, 20 marks): ``Colonial administrators helped to construct the very traditionalism which marked the Indian society as backward. Comment.''}
\item \textbf{Introduction}: 'the above statement reflects the reaction to the \textbf{Indophobic sentiments} which had downgraded India as backward.'
\end{itemize}
\end{PYQLink}

\begin{KeyConcept}
\begin{itemize}
\item \textbf{Bernard Cohn}: '\textbf{history finds data and social anthropology creates data}' --- and his student \textbf{Nicholas Dirks} ('Castes of Mind'): '\textbf{colonialism not only creates knowledge but is produced by it}'.
\item \textbf{Bernard Cohn's key view}: the British conquered not only the territory of India but also its \textbf{epistemological space} --- i.e. \emph{how we know India} (the ancient/medieval/modern frame) is itself a colonial creation. This 'cultural analysis' was labelled \textbf{ethno-sociology} (McKim Marriott's term) --- creating knowledge about the people you will rule, to justify ruling them.
\item Cohn: the British devised \textbf{investigative modalities} --- \textbf{historiography, observation and travel, survey/enumeration (census), museology and surveillance} --- to collect and generate knowledge of the dominated world; the \textbf{accumulation of this knowledge, along with the exercise of power, were parts of the same project of colonialism}.
\end{itemize}
\end{KeyConcept}

\begin{KeyConcept}
\begin{itemize}
\item The instruments, one by one:
\item \textbf{Census} --- \textbf{Ghurye's classical statement}: 'through the census, the British tried to create \textbf{solutions for the problems created by them only}.' Caste as a \textbf{pan-India phenomenon is a British creation}, and through the census they wanted to give quota/representation to different caste groups --- solving a problem of their own making. (The example is \textbf{L. Middleton}, Punjab census commissioner, 1921.) The census produced \textbf{caste patriotism, caste associations and caste divisions}.
\item \textbf{Historiography} --- history is not only about finding data but \textbf{creating data}: colonial historiography shows village India as static, backward, self-sufficient and a 'little republic'.
\item \textbf{Observation and travel} --- British travellers and missionaries (e.g. \textbf{Francis Buchanan}, the Jesuits) went into kingdoms, observed and studied their civilisations, and thereby 'created knowledge' about India as a whole.
\item \textbf{Museum} --- a museum lets an outsider 'know' the inside, presenting India in a chronological/evolutionary way, just as the British wanted to know India from its (constructed) past.
\item \textbf{Gazetteers} (and museums) are equally instruments of this knowledge-creation --- to be remembered alongside census, historiography and observation/travel.
\item \textbf{Contemporary 'colonial hangover'} --- the colonialists are gone, but the \textbf{caste census still runs}, and it still shows Indian society as backward. Even though the caste census has advantages (delineating policy, channelling resources to unrepresented sections), its \textbf{latent dysfunction is that it recreates caste consciousness} in a free, sovereign, equality-speaking India --- which is why \textbf{Ambedkar's} warning that \textbf{Hinduism and democracy cannot survive together} is now more relevant than ever.
\end{itemize}
\end{KeyConcept}

\begin{QuickRevision}
\begin{itemize}
\item PYQ 2020 (20 marks): 'colonial administrators helped construct the very traditionalism that marked India as backward.'
\item Intro: reaction to Indophobic sentiments; Bernard Cohn ('history finds data, social anthropology creates data') + Nicholas Dirks ('colonialism creates knowledge and is produced by it').
\item Cohn: British conquered the 'epistemological space'; investigative modalities (historiography, observation/travel, survey/enumeration, museology, surveillance); ethno-sociology (Marriott).
\item Instruments: census (Ghurye's 'solution for problems created by them'; L. Middleton 1921; caste as pan-India creation; caste patriotism), historiography (creating data), observation/travel (Buchanan, Jesuits), museum (evolutionary view).
\item Contemporary: colonial hangover --- caste census still runs; latent dysfunction = recreates caste consciousness; Ambedkar (Hinduism vs democracy).
\end{itemize}
\end{QuickRevision}

\section{PYQ: Is Indian Rural Society a Faction-Ridden Society?}

\begin{PYQLink}
\begin{itemize}
\item \textbf{PYQ: ``Indian rural society is a faction-ridden society. Discuss.''} A 'faction' = a division.
\item The scholars' field studies all reveal \textbf{multiple divisions within rural society}: \textbf{Srinivas} (Rampura --- Vokkaligas, Lingayats, and the 'untouchable' smiths), \textbf{S. C. Dube} (Shamirpet --- Hindu jatis, Muslim jatis, untouchables), \textbf{Beteille} (Shripuram --- Brahmins, non-Brahmins, Adi Dravidas), and \textbf{Moffatt} (Andavur --- who adds one more dimension: the '\textbf{Brahmins of the untouchables}').
\end{itemize}
\end{PYQLink}

\begin{KeyConcept}
\textbf{Conclusion}: seen through the \textbf{colonial/Indological lens}, India consists of homogeneous elements (the village, the Brahmins); but seen through the \textbf{field-worker's lens}, India consists of \textbf{multiple Brahmins in multiple villages} --- finding multiple Brahmins within a single village vindicates the claim that Indian rural society is \textbf{faction-ridden} (there is not one Brahmin but many factions within the Brahmins).
\end{KeyConcept}

\begin{QuickRevision}
\begin{itemize}
\item Faction = division; rural society has multiple divisions.
\item Srinivas (Rampura), Dube (Shamirpet), Beteille (Shripuram), Moffatt (Andavur 'Brahmins of the untouchables').
\item Colonial/Indological lens = homogeneous India; field lens = multiple Brahmins in multiple villages $\rightarrow$ faction-ridden.
\end{itemize}
\end{QuickRevision}

